Mathematical National Park has a collection of trails. There are designated
campsites along the trails, including a campsite at each intersection of
trails. The rangers call each stretch of trail between adjacent campsites a
``segment". The trails have been laid out so that it is possible to take a hike
that starts at any campsite, covers each segment exactly once, and ends at the
beginning campsite. Prove that it is possible to plan a collection
$\mathcal{C}$ of hikes with all of the following properties:

\begin{enumerate}[(a)]

\item Each segment is covered exactly once in one hike $h \in \mathcal{C}$ and
      never in any of the other hikes of $\mathcal{C}$.
\item Each $h \in \mathcal{C}$ has a base campsite that is its beginning and
      end, but which is never passed in the middle of the hike. (Different hikes of
      $\mathcal{C}$ may have different base campsites.)
\item Except for its base campsite at beginning and end, no hike in
      $\mathcal{C}$ passes any campsite more than once.

\end{enumerate}
